\scanpage{215}\mspage{159}
% archive sheet number: 214

% MS 1–8
Mais je m'aperçois à l'instant qu'on \unsure{trouve|tourne} quasiment ``à l'œil'' le fait que $f_0^*$ soit bien une \uline{équivalence} de \unsure{catégorie|catégories}, et que $f_!^0$ soit l'équivalence quasi-inverse, en utilisant l'adjonction des foncteurs $f_!^0,f_0^*$, et se \unsure{rappellant|rappelant} que pour prouver que ce sont des équivalences quasi-inv. l'une de l'autre, il \deleted{suffit}\inserted{revient au même} de montrer que les morphismes d'adjonction

% MS 9–10
\[
\begin{aligned}
 \mathrm{id}&\xrightarrow{\alpha_0} f_0^* f_!^0\\
 \mathrm{id}&\xleftarrow{\beta_0} f_!^0 f_0^*
\end{aligned}
\]

% MS 11
sont \uline{iso}.

% MS 11–14 crossed block
\deleted{J'admets que ces flèches sont données par le $=$ \unsure{procédé|précédent} que dans le cas de $f_!,f^*$, 1) $\alpha_0$ est iso. Je prends $\unsure{\xi|\zeta}'\in\mathrm{Ob}\,\mathrm{HOT}^{\mathrm{lc}}_{\unsure{W|w}}(S')$, il est défini par $X'\in\mathrm{Ob}\,\deleted{\mathrm{Cat}_{/S'}}\,\unsure{\mathscr P|\mathscr S}_{S'}$ comme}

% MS 15–16 crossed display
\[
\deleted{\begin{aligned}
 \xi'&=\mathrm{hot}_{S',W}^{\mathrm{lc}}(X'),\ \text{d'où}\quad
 f_!^0(\xi')=\mathrm{hot}_{S,W}^{\mathrm{lc}}(X'),\quad
 f_0^*f_!^0(\xi')=\mathrm{hot}_{S',W}(C_S(X')\times_S S'),\\
 &\hspace{8em}=\mathrm{hot}_{S,W}(C(\unsure{X'|X}))
\end{aligned}}
\]

% MS 17–19
Je vais essayer de décrire ces flèches d'adjonction dans le cas général (sans supposer $f\in W$), puis de montrer qu'elles sont dans $W$ si $f$ l'est.

% MS 20
$\xi=\mathrm{hot}_{S,W}^{\mathrm{l.c.}}(X)$ dans $\mathrm{HOT}_W(S)$ \unsure{OTS|OOT|OOI} $X\in\mathrm{Ob}\,\uline{\mathrm{Parf}}\,S$.

% MS 21
\deleted{Comme} $X$ lisse sur $S$, on a

% MS 22–23
\[
\begin{aligned}
 f_0^*(\xi)&=\mathrm{hot}_{S',W}^{\mathrm{lc}}(X\times_S S'),\quad
 f_!^0 f_0^*(\xi)=\mathrm{hot}_{S,W}^{\mathrm{lc}}(X\times_S S')\\
 &\hspace{12em}\deleted{{}=\mathrm{hot}_{S,W}^{\mathrm{lc}}(C_S(X\times_S S'))}
\end{aligned}
\]

% MS 24–25
Donc \deleted{\unsure{$f_0^*$})} $\beta_{0\,\xi}:f_!^0f_0^*(\xi)\longrightarrow\xi$ flèche

% MS 26
\[
 X\times_S S'\longrightarrow X\qquad\text{dans }(\mathrm{Cat}_{/S})(W_S)^{-1}
\]

% MS lower left, alongside 24–27
% slanted
\margin{\uline{NB} On voit que l'axiome L6 est $\pm$ \uline{nécessaire} pour que la \unsure{validité} du \unsure{calcul} $f_0^*$, $f_!^0$ \unsure{soit|ait une} \unsure{valable|valeur}}

% MS 27–28
Je présume que c'est \unsure{$\mathrm{pr}_1$|$\mathrm{pr}\,1$}. Est-ce dans $W$ si $f:S'\longrightarrow S$ l'est ? Oui, car $X$ est parfait
